\begin{lemma}[Jensen's lemma]
Let $E$ be a finite-dimensional real inner-product space with an additive Haar
measure $\mu$, let $\lambda>0$ and $r>0$, let $z_0\in E$ and let
$\phi:E\to\R$ be continuous on $\overline B(z_0,r)$ and $\lambda$-semiconvex on
$B(z_0,r)$ in the sense of \noderef{SemiconvexOn}. Assume that $z_0$ is a strict
maximum point of $\phi$ on the closed ball:
\[
  \phi(z)<\phi(z_0)\qquad\bigl(z\in\overline B(z_0,r),\ z\ne z_0\bigr).
\]
Then for every $\delta>0$
\[
  \mu\Bigl(\bigl\{z\in B(z_0,r)\ :\ \exists\,p\in E,\ |p|\le\delta,\
    w\mapsto\phi(w)+\langle p,w\rangle\ \text{has a local maximum on }
    B(z_0,r)\text{ at }z\bigr\}\Bigr)>0 .
\]
\end{lemma}
\begin{proof}
Write $B=B(z_0,r)$ and fix the three radii
\[
  r_1=\frac r2,\qquad \rho=\frac r8,\qquad \eta=\frac{3r}{8},
\]
so that $0<\rho\le r_1<r$, $\eta>0$ and $\rho+\eta=r_1$. Write
\[
  \psi(w)=\phi(w)+\frac{\lambda|w|^2}{2}\qquad(w\in E),
\]
which is convex on the open convex set $B$ by the definition of
$\lambda$-semiconvexity in \noderef{SemiconvexOn}. Denote by $K_\delta$ the set
appearing in the statement. Finally put $n=\dim E$.

No nondegeneracy assumption on $E$ is made or needed: the only place where the
old argument used a nonzero vector was the construction of the threshold in
Step 1, and that construction has been moved into
\noderef{StrictMaxPerturbLocalization}, which covers the degenerate case
internally. The case $n=0$ is discussed again at the end of Step 5.

\emph{Step 1: the localization threshold.}
Since $r_1<r$ we have $\overline B(z_0,r_1)\subseteq\overline B(z_0,r)$, so
$\phi$ is continuous on $\overline B(z_0,r_1)$ and
\[
  \phi(z)<\phi(z_0)\qquad\bigl(z\in\overline B(z_0,r_1),\ z\ne z_0\bigr),
\]
both being restrictions of the corresponding hypotheses of the lemma. Together
with $0<\rho\le r_1$ these are exactly the hypotheses of
\noderef{StrictMaxPerturbLocalization} for $\phi$, $z_0$, $r_1$, $\rho$. That
lemma therefore supplies $\delta_0>0$ such that
\begin{equation}\label{eq:jl-loc}
  \begin{gathered}
  \text{if }|p|\le\delta_0,\ z\in\overline B(z_0,r_1)\text{ and }
  \phi(w)+\langle p,w\rangle\le\phi(z)+\langle p,z\rangle
  \text{ for all }w\in\overline B(z_0,r_1),\\
  \text{then }|z-z_0|<\rho .
  \end{gathered}
\end{equation}
Put
\[
  d=\min\{\delta,\delta_0\}>0 ,
\]
so that $d\le\delta$ and $d\le\delta_0$.

\emph{Step 2: the set $K$ and its elementary properties.}
Define
\[
  K=\Bigl\{z\in\overline B(z_0,r_1)\ :\ \exists\,p\in E,\ |p|\le d,\
    \phi(w)+\langle p,w\rangle\le\phi(z)+\langle p,z\rangle\
    \text{for all}\ w\in\overline B(z_0,r_1)\Bigr\}.
\]
Call such a $p$ a \emph{multiplier for $z$}.

\emph{(2a) $K\subseteq B(z_0,\rho)$.} If $z\in K$ with multiplier $p$, then
$|p|\le d\le\delta_0$, so \eqref{eq:jl-loc} applies and gives $|z-z_0|<\rho$.

\emph{(2b) Room around a localized point.} If $|z-z_0|<\rho$ then
\[
  B(z,\eta)\subseteq\overline B(z_0,r_1)\cap B .
\]
Indeed, for $w\in B(z,\eta)$ the triangle inequality gives
\[
  |w-z_0|\le|w-z|+|z-z_0|<\eta+\rho=\frac{3r}{8}+\frac r8=\frac r2=r_1,
\]
so $|w-z_0|\le r_1$, i.e.\ $w\in\overline B(z_0,r_1)$, and also $|w-z_0|<r_1<r$,
i.e.\ $w\in B$.

\emph{(2c) $K\subseteq K_\delta$.} Let $z\in K$ with multiplier $p$. By (2a)
$|z-z_0|<\rho$, so (2b) applies to $z$; in particular
$z\in B(z,\eta)\subseteq B$. The set $B(z,\eta)$ is open and contains $z$,
hence is a neighbourhood of $z$, and $B(z,\eta)\subseteq\overline B(z_0,r_1)$ by
(2b); therefore
\[
  \phi(w)+\langle p,w\rangle\le\phi(z)+\langle p,z\rangle
  \qquad\bigl(w\in B(z,\eta)\bigr),
\]
and a fortiori for $w\in B(z,\eta)\cap B$. Thus
$w\mapsto\phi(w)+\langle p,w\rangle$ has a local maximum on $B$ at $z$. Since
also $|p|\le d\le\delta$, the point $z$ lies in $K_\delta$.

\emph{(2d) Every admissible multiplier occurs.} Let $p\in E$ with $|p|\le d$.
The set $\overline B(z_0,r_1)$ is closed and bounded, hence compact (preamble),
and it is nonempty since $z_0$ belongs to it ($r_1>0$). The function
$w\mapsto\phi(w)+\langle p,w\rangle$ is continuous on it, being the sum of
$\phi$ --- continuous on $\overline B(z_0,r_1)\subseteq\overline B(z_0,r)$ ---
and of the continuous linear functional $w\mapsto\langle p,w\rangle$. By
Weierstrass' theorem (preamble) it attains a maximum at some
$z\in\overline B(z_0,r_1)$, i.e.
\[
  \phi(w)+\langle p,w\rangle\le\phi(z)+\langle p,z\rangle
  \qquad\bigl(w\in\overline B(z_0,r_1)\bigr).
\]
Then $z\in K$ and $p$ is a multiplier for $z$.

\emph{Step 3: the multiplier is unique, so $H$ is well defined.}
Let $z\in K$ and let $p,p'$ both be multipliers for $z$. By (2a) and (2b),
$B(z,\eta)\subseteq\overline B(z_0,r_1)\cap B$; in particular $z\in B$. Since
$B$ is open and $\psi$ is convex on $B$, \noderef{ConvexSubgradientExists}
applied to $B$, $\psi$ and $z$ produces $s\in E$ with
\begin{equation}\label{eq:jl-sub}
  \psi(z)+\langle s,w-z\rangle\le\psi(w)\qquad(w\in B).
\end{equation}
Apply \noderef{MaxMultiplierPinsSubgradient} with the data
\[
  \lambda,\quad\phi,\quad A=\overline B(z_0,r_1),\quad\Omega=B,\quad
  z,\quad p,\quad s,\quad\eta .
\]
Its three hypotheses hold: $\eta>0$ and
$B(z,\eta)\subseteq\overline B(z_0,r_1)\cap B$ as just noted; the maximality of
$\phi+\langle p,\cdot\rangle$ at $z$ over $A$, because $p$ is a multiplier for
$z$; and \eqref{eq:jl-sub}, which is the subgradient inequality for $\psi$ on
$\Omega=B$. The conclusion is $s=\lambda z-p$. The same lemma with $p'$ in place
of $p$ --- the subgradient $s$ being unchanged, since it was produced from $z$
alone --- gives $s=\lambda z-p'$. Hence $\lambda z-p=\lambda z-p'$ and so
$p=p'$.

Therefore each $z\in K$ has exactly one multiplier, which we call $H(z)$; by
definition $|H(z)|\le d$ and
\begin{equation}\label{eq:jl-mult}
  \phi(w)+\langle H(z),w\rangle\le\phi(z)+\langle H(z),z\rangle
  \qquad\bigl(z\in K,\ w\in\overline B(z_0,r_1)\bigr).
\end{equation}
Extend $H$ to all of $E$ by setting $H(z)=0$ for $z\notin K$; this does not
affect $H(K)$ or the Lipschitz estimate of Step 4, both of which involve only
points of $K$.

\emph{Step 4: $H$ is $2\lambda$-Lipschitz on $K$.}
Let $z_1,z_2\in K$. We verify the hypotheses of
\noderef{SemiconvexMultiplierLipschitz} for the data
\[
  \lambda,\quad\Omega=B,\quad\phi,\quad A=\overline B(z_0,r_1),\quad
  z_1,\ z_2,\quad p_1=H(z_1),\ p_2=H(z_2),\quad\eta .
\]
First, $\lambda>0$ by hypothesis, $B$ is open, and $\phi$ is
$\lambda$-semiconvex on $B$; and $\eta>0$. Next, by (2a) each $z_i$ satisfies
$|z_i-z_0|<\rho$, so (2b) gives
$B(z_i,\eta)\subseteq\overline B(z_0,r_1)\cap B$ for $i=1,2$, which is
hypothesis (1) of that lemma. For hypothesis (2),
\[
  |z_1-z_2|\le|z_1-z_0|+|z_0-z_2|<\rho+\rho=\frac r4<\frac{3r}{8}=\eta .
\]
Hypothesis (3) is \eqref{eq:jl-mult} for $z_1$ and for $z_2$. The conclusion is
\begin{equation}\label{eq:jl-lip}
  |H(z_1)-H(z_2)|\le2\lambda|z_1-z_2|\qquad(z_1,z_2\in K).
\end{equation}

\emph{Step 5: the measure estimate.}
By \eqref{eq:jl-lip} the map $H:E\to E$ satisfies
$|H(a)-H(a')|\le2\lambda|a-a'|$ for all $a,a'\in K$, and $2\lambda>0$. So
\noderef{LipschitzImageHaarBound}, applied with $A=K$, $f=H$ and the constant
$2\lambda$, gives
\begin{equation}\label{eq:jl-image}
  \mu\bigl(H(K)\bigr)\le(2\lambda)^{n}\,\mu(K).
\end{equation}
(No measurability of $K$ is needed: that lemma has none, and below we use only
the monotonicity of $\mu$, which holds for arbitrary sets.)

Next, $\overline B(0,d)\subseteq H(K)$. Indeed, let $|p|\le d$. By (2d) there is
$z\in K$ for which $p$ is a multiplier, and by Step 3 the multiplier of $z$ is
unique and equals $H(z)$; hence $p=H(z)\in H(K)$.

By the preamble's computation of the Haar measure of balls,
\[
  \mu\bigl(\overline B(0,d)\bigr)=d^{\,n}\,\mu\bigl(B(0,1)\bigr)
  \qquad\text{with}\qquad 0<\mu\bigl(B(0,1)\bigr)<\infty,
\]
and $d>0$, so $\mu(\overline B(0,d))$ is a strictly positive (finite) number.
Combining this with the monotonicity of $\mu$ along
$\overline B(0,d)\subseteq H(K)$ and with \eqref{eq:jl-image},
\[
  0<\mu\bigl(\overline B(0,d)\bigr)\le\mu\bigl(H(K)\bigr)
   \le(2\lambda)^{n}\,\mu(K).
\]
The factor $(2\lambda)^{n}$ is a finite real number, so if $\mu(K)$ were $0$ the
right-hand side would be $0$, contradicting the strict positivity of the
left-hand side. Hence
\[
  \mu(K)>0 .
\]
This argument is insensitive to the dimension: when $n=0$ both $d^{\,n}$ and
$(2\lambda)^{n}$ equal $1$, the chain reads
$0<\mu(B(0,1))\le\mu(H(K))\le\mu(K)$, and the conclusion $\mu(K)>0$ is
unchanged.

Finally $K\subseteq K_\delta$ by (2c), so the monotonicity of $\mu$ gives
\[
  \mu(K_\delta)\ge\mu(K)>0 ,
\]
which is the assertion of the lemma.
\end{proof}
